\section{Objetivos Específicos}

Os objetivos específicos deste trabalho incluem:

\begin{itemize}
    \item Desenvolver um \textit{pipeline} reprodutível para treino e avaliação de modelos de redução de artefatos em\acrshort{ct} no domínio da imagem, com separação rigorosa entre dados de treino, validação e teste.

    \item Desenvolver uma arquitetura \textit{baseline} baseada em\acrshort{unet} para correção de artefatos em imagens de\acrshort{ct}, garantindo reprodutibilidade experimental.

    \item Avaliar uma arquitetura\acrshort{unet} \textit{single-head} para reconstrução de imagem como referência metodológica.

    \item Propor e avaliar uma arquitetura\acrshort{unet} \textit{dual-head} baseada em aprendizado multi-tarefa para reconstrução de imagem e predição de artefatos, incluindo um estudo de ablação comparativo entre as duas abordagens.

    \item Avaliar o impacto de módulos arquiteturais avançados, incluindo\textit{\acrfull{aspp}},\textit{\acrfull{se}} e\textit{\acrfull{cbam}}, por meio de estudos de ablation controlados.

    \item Definir e aplicar critérios objetivos de seleção do modelo com base em consistência entre conjuntos de validação e teste, análise de distribuições de desempenho, controlo de sobreajuste e testes estatísticos pareados.

    \item Avaliar o desempenho dos modelos utilizando métricas quantitativas (\acrshort{psnr},\acrshort{ms-ssim},\acrshort{ssim},\acrshort{bce} e\acrshort{dice} ), análise qualitativa e mapas de erro.

    \item Investigar a capacidade de generalização do modelo selecionado por meio de validação externa em dados independentes provenientes do\textit{\acrfull{tcia}}.
\end{itemize}